
\documentclass{article}
\usepackage{colortbl}
\usepackage{makecell}
\usepackage{multirow}
\usepackage{supertabular}

\begin{document}

\newcounter{utterance}

\centering \large Interaction Transcript for game `cladder', experiment `full\_v1.5\_default', episode 3456 with selfrefine{-}Qwen3.5{-}9B.
\vspace{24pt}

{ \footnotesize  \setcounter{utterance}{1}
\setlength{\tabcolsep}{0pt}
\begin{supertabular}{c@{$\;$}|p{.15\linewidth}@{}p{.15\linewidth}p{.15\linewidth}p{.15\linewidth}p{.15\linewidth}p{.15\linewidth}}
    \# & \multicolumn{2}{c}{Player} && \multicolumn{2}{c}{Game Master} \\
    \hline

    \theutterance \stepcounter{utterance}  
    & & & \multicolumn{4}{p{0.6\linewidth}}{
        \cellcolor[rgb]{0.9,0.9,0.9}{
            \makecell[{{p{\linewidth}}}]{
                \texttt{\tiny{[P1$\langle$GM]}}
                \texttt{You are an expert in causal inference. The following question is not a typical commonsense query, but rather a meticulously designed question created by a professor specializing in causal inference, intended to assess the students' mastery of the course content.} \\
\texttt{Imagine a self{-}contained, hypothetical world with only the following conditions, and without any unmentioned factors or causal relationships: Solar eclipse has a direct effect on Alice arriving to school. Traffic has a direct effect on Alice arriving to school. Method 1: We look at how solar eclipse correlates with Alice arriving to school case by case according to traffic. Method 2: We look directly at how solar eclipse correlates with Alice arriving to school in general. To understand how solar eclipse affects Alice arriving to school, is it more correct to use the Method 1 than Method 2?} \\
\texttt{Start your answer with 'Yes' or 'No', followed by additional reasoning or evidence to support your explanation.} \\
            }
        }
    }
     \\ \\

    \theutterance \stepcounter{utterance}  
    & & \multicolumn{4}{p{0.6\linewidth}}{
        \cellcolor[rgb]{0.95,0.95,0.95}{
            \makecell[{{p{\linewidth}}}]{
                \texttt{\tiny{[P1$\leftrightarrow$P]}}
                \texttt{{[}DRAFT{]}} \\
\texttt{No.} \\
\\ 
\texttt{In this hypothetical world, **Traffic** acts as a **confounder** because it has a direct effect on the outcome (Alice arriving to school) and, by the structure of the question's logic, is the variable that must be controlled to isolate the specific causal mechanism of the solar eclipse. If we use **Method 2** (looking at the general correlation), we are observing the **marginal association** between the solar eclipse and arrival time. This marginal association will be biased because it aggregates data across different levels of traffic. Depending on the distribution of traffic conditions, the observed correlation could be diluted, exaggerated, or even reversed (M{-}bias or Simpson's paradox scenarios), failing to reveal the true causal effect of the eclipse.} \\
\\ 
\texttt{**Method 1** describes the process of **stratification** (or conditioning on the confounder). By looking at the correlation "case by case according to traffic," we are effectively calculating the conditional probability \$P(\textbackslash{}text\{Arrival\} | \textbackslash{}text\{Eclipse\}, \textbackslash{}text\{Traffic\})\$. In causal inference, to identify the causal effect of a treatment (Solar Eclipse) on an outcome (Arrival), one must block all back{-}door paths. Since Traffic is a common cause (or a variable influencing the outcome that varies independently of the eclipse's specific mechanism in this context), failing to condition on it leaves the estimate biased. Therefore, Method 1 is the correct approach to isolate the direct causal effect, whereas Method 2 yields a confounded estimate.} \\
            }
        }
    }
    & \\ \\

    \theutterance \stepcounter{utterance}  
    & & \multicolumn{4}{p{0.6\linewidth}}{
        \cellcolor[rgb]{0.95,0.95,0.95}{
            \makecell[{{p{\linewidth}}}]{
                \texttt{\tiny{[P1$\leftrightarrow$P]}}
                \texttt{{[}SELF{-}FEEDBACK{]}} \\
\texttt{FORMAT: ok} \\
\texttt{RULES: ok} \\
\texttt{STRATEGY: problem} \\
\texttt{REASON: The answer incorrectly identifies Traffic as a confounder; since the prompt states Traffic has a direct effect on the outcome but does not state it causes the Solar Eclipse, Traffic is a mediator or independent cause, not a confounder, making the causal reasoning flawed.} \\
            }
        }
    }
    & \\ \\

    \theutterance \stepcounter{utterance}  
    & \multicolumn{4}{p{0.6\linewidth}}{
        \cellcolor[rgb]{0.561,0.737,0.561}{
            \makecell[{{p{\linewidth}}}]{
                \texttt{\tiny{[P1$\rangle$GM]}}
                \texttt{No.} \\
\\ 
\texttt{In this hypothetical world, **Traffic** is **not a confounder** because the prompt explicitly states it has a direct effect on the outcome but does not state that it causes the Solar Eclipse. A confounder must be a common cause of both the treatment (Solar Eclipse) and the outcome (Arrival). Since Traffic does not cause the Eclipse, it cannot be a confounder in the traditional sense.} \\
\\ 
\texttt{Instead, the relationship between the Solar Eclipse and Traffic is undefined in the prompt. If the Solar Eclipse and Traffic are independent events, then the marginal correlation (Method 2) is a valid estimate of the causal effect, as there is no back{-}door path to block. Method 1 (stratification) would only be necessary if Traffic were a **mediator** (if the Eclipse caused Traffic, which caused Arrival) or if there were an unmentioned common cause. However, based strictly on the provided text, assuming Traffic is a confounder is a logical error. Therefore, Method 1 is not necessarily more correct than Method 2; in fact, stratifying on an independent variable can sometimes introduce unnecessary variance or misinterpretation if the goal is simply to estimate the total effect. The correct approach depends on the unknown relationship between the Eclipse and Traffic, not on a false assumption of confounding.} \\
            }
        }
    }
    & & \\ \\

    \theutterance \stepcounter{utterance}  
    & & & \multicolumn{2}{p{0.3\linewidth}}{
        \cellcolor[rgb]{0.9,0.9,0.9}{
            \makecell[{{p{\linewidth}}}]{
                \texttt{\tiny{[GM$|$GM]}}
                \texttt{No} \\
            }
        }
    }
    & & \\ \\

    \theutterance \stepcounter{utterance}  
    & & & \multicolumn{2}{p{0.3\linewidth}}{
        \cellcolor[rgb]{0.9,0.9,0.9}{
            \makecell[{{p{\linewidth}}}]{
                \texttt{\tiny{[GM$|$GM]}}
                \texttt{no} \\
            }
        }
    }
    & & \\ \\

    \theutterance \stepcounter{utterance}  
    & & & \multicolumn{2}{p{0.3\linewidth}}{
        \cellcolor[rgb]{0.9,0.9,0.9}{
            \makecell[{{p{\linewidth}}}]{
                \texttt{\tiny{[GM$|$GM]}}
                \texttt{game\_result = WIN} \\
            }
        }
    }
    & & \\ \\

\end{supertabular}
}

\end{document}
